\chapter{Host organisation and research context}
\label{ch:host}

\section{Forschungszentrum Jülich and IAS-7}

The internship was carried out at Forschungszentrum Jülich, a member of the
Helmholtz Association, at IAS-7 (Civil Safety Research), part of the
Institute for Advanced Simulation and directed by Prof.\ Dr.\ Armin Seyfried.
IAS-7 works on pedestrian and fire dynamics through combined experimental and
simulation approaches. The link between an insect registration pipeline and
an institute concerned with crowd safety runs through method, not subject
matter: IAS-7's long-running tool \term{PeTrack} recovers pedestrian
trajectories from video automatically, because manual extraction cannot keep
pace with the institute's experimental series, but a trajectory is only a
point per person per frame. \term{smilify} is designed to give the same kind
of automated recovery a richer target: full three-dimensional pose and
shape, species-agnostic, from monocular or multi-view RGB video, and ants
serve as a demanding test case for it: thinner structures, more limbs,
greater articulation, and worse input data than a human subject in a
laboratory corridor. A registration method proven on a six-legged,
few-millimetre specimen from imperfect photogrammetry is tested well beyond
the conditions of its eventual deployment.

Shared allocation on the Jülich Supercomputing Centre's own machines could
not absorb the internship's full compute demand; the more resource-intensive
experiments in this report were instead run on external HPC infrastructure
(\S\ref{sec:infrastructure}), so the computational scale of the work was not
a direct consequence of being hosted at Jülich.

\section{Research context: parametric models beyond human subjects}

This internship extends parametric three-dimensional shape models to species
outside the small set computer vision has historically addressed.
\term{SMPL}~\cite{smpl} assumes a cooperative subject, a large scan corpus,
and a fixed body topology; \term{SMAL}~\cite{smal} relaxed the first two
assumptions for quadrupeds by learning a shape space from toy figurines.
\term{smilify} continues this progression to insects, where the body plan
itself differs and the structures of interest are thin, numerous, and
highly articulated (Chapter~\ref{ch:pipeline} traces this lineage in full).
Upstream sits \term{scAnt}~\cite{scant}, an open-source macro photogrammetry
scanner by the same author, which makes high-resolution 3D capture of small
specimens accessible without dedicated scanning infrastructure: an
affordable scanner producing abundant imperfect data, and a robust
registration pipeline converting it into comparable measurements; neither
has much independent value without the other, and the failure mode
identified in Chapter~\ref{ch:introduction} sits exactly at their junction.

\section{Supervision and working arrangement}

The internship was supervised by Dr.\ Fabian Plum, postdoctoral researcher
at IAS-7 and author of \term{scAnt}~\cite{scant}, \term{replicAnt}~\cite{replicant},
\term{OmniTrax}~\cite{omnitrax}, and \term{smilify}, and ran on site from 1
July to 25 September 2026. Two aspects of that arrangement shaped the work.
First, the technical data sheet's assignments were an initial hypothesis
about where effort would be valuable, not a fixed specification: when the
controlled comparison in \S\ref{sec:turn} showed mesh input quality was not
the limiting factor, redirecting the remaining work toward initialisation
and correspondence was the correct response to the evidence. Second,
development took place on a live, shared repository under standard
pull-request review, with the supervisor running related experiments on
adjacent branches; several results reported later in this document were
revised following that review, and the pre-registration discipline used
throughout the report follows directly from it. Beyond ants, the same
pipeline has since been used to build parametric models of other species,
context worth noting since the problems investigated here are not specific
to Formicidae.

\section{Academic framing}

This internship forms part of the second year of the engineering cycle at
Universit\'e Mohammed VI Polytechnique, and this report is the corresponding
academic deliverable. It documents both the technical contributions produced
and the evidence behind the decisions that shaped the work's direction,
inseparable here, since the most consequential decisions taken during the
internship concerned which lines of investigation to discontinue, and those
decisions are only defensible alongside the evidence that motivated them.
